\section{Artificial intelligence as several different shocks}\label{sec:ai}

``AI'' enters the model through several parameters, and each leaves a different
footprint. That is what makes the model testable rather than a story.

\begin{proposition}[Signatures of AI]\label{prop:ai}
In the large-topic limit:
\begin{enumerate}[label=(\roman*)]
\item \emph{Faster drafting} ($a\uparrow$) acts like a larger community. Density moves with elasticity $\mathcal E$ (zero if $\eta=0$) and the number of topics with elasticity $1-D\mathcal E$. Papers per researcher rise with elasticity $1+\eps_x\mathcal E$, and output per researcher with elasticity $1-\mathcal E|\eps_w|$.
\item \emph{Cheaper entry into science} ($c\downarrow$, with $N$ set by $aF(M)=\omega+c$ and $\eta>0$) raises $N$. Density rises with elasticity $\mathcal E$ with respect to $N$, papers per researcher move with elasticity $\eps_x\mathcal E\approx 0$, and new topics absorb a share $1-D\mathcal E$ of growth.
\item \emph{Better reading, learning or search, or cheaper pivots} ($b$, $\beta$ or $\sigma\uparrow$, $\kappa_p\downarrow$) raise $v$ and $w$ at every density. Density rises and the number of topics falls: science concentrates. Output per researcher rises only through the frontier, by $\eta\,|\dd\ln M^*|$, and is unchanged when $\eta=0$.
\item \emph{Faster catch-up} ($b_c\uparrow$, e.g.\ literature-summarising tools) raises $w$ at every density, with the effects in (iii). It also lowers the overshooting threshold $g^*$ of \cref{prop:toofast}, so more topics grow too fast.
\item \emph{Better search} ($\sigma\uparrow$) raises the aware share $A$ at any given reading coverage.
\end{enumerate}
\end{proposition}

Parts (iii) and (iv) follow from the envelope theorem. Each parameter raises
$g(r,P)$ at every $(r,P)$, hence raises $v$. A pointwise upward shift of a
single-peaked curve moves its falling-side crossing with a fixed $F$ to the right.
The effect on the measured overlap rate is ambiguous in general: denser topics
generate more overlap, but better detection removes more of it.
\Cref{tab:ai} reports the calibrated signs.

\begin{table}[t]
\centering
\caption{The footprints of AI in the calibrated model}\label{tab:ai}
\small
\begin{tabular}{@{}llrrrrrrr@{}}
\toprule
Channel & Shock & Density & Topics & Papers/ & Output/ & Flagged & Aware & $g^*$ \\
 & & & & researcher & researcher & overlap & share (pp) & (pp) \\
\midrule
Faster drafting & $a\times1.3$ & +5.1 & +20.0 & +26.1 & +25.4 & +6.7 & -1.3 & +0.0 \\
Cheaper entry & $N\times1.3$ & +5.1 & +20.0 & -0.1 & -0.8 & +6.7 & -1.3 & +0.0 \\
Faster reading & $b\times1.3$ & +8.5 & -9.9 & +0.2 & +0.4 & +1.6 & +5.8 & +0.5 \\
Better learning & $\beta\times1.3$ & +13.6 & -16.9 & -0.5 & +0.7 & +17.0 & -2.7 & -5.6 \\
Better search & $\sigma\times1.3$ & +0.3 & -0.4 & +0.0 & +0.0 & -12.4 & +11.1 & +0.6 \\
Faster catch-up & $b_c\times 2$ & +61.3 & -61.3 & -1.5 & +2.6 & +79.5 & -13.6 & -23.6 \\
Cheaper pivots & $\kappa_p\div 2$ & +1.1 & -1.4 & -0.1 & +0.1 & -21.6 & -15.3 & +0.6 \\
\bottomrule
\end{tabular}
\par\smallskip\footnotesize Long-run stationary effects in log points ($\times100$) of a 30\% shock (catch-up speed doubled, pivot cost halved). Density is $P^*$; output per researcher is $aF(M^*)$. The flagged overlap rate is relative to its baseline level, which is a normalisation (\cref{sec:facts}). $g^*$ is the overshooting threshold of \cref{prop:toofast}.

\end{table}

\paragraph{What the arXiv data say.} The footprints differ enough to tell the
channels apart. Four facts settle the question for 2023--2026.
\begin{itemize}
\item Papers per author stayed flat while authors grew (Fact~\ref{fact:authors}). Faster
drafting raises papers per author one for one. Cheaper entry leaves them flat.
\item High-uptake areas produced \emph{more} pioneer papers and \emph{lower} local
density relative to their own pre-ChatGPT trend (Fact~\ref{fact:concentration}). All reading-side channels predict
concentration: fewer new topics and denser old ones. Cheaper entry predicts more
new topics. Faster growth also raises the share of young, thin topics in the
life-cycle model, which lowers mean density.
\item The aware share of close overlap stayed at 35--42\% (Fact~\ref{fact:aware}). Better search
would have raised it.
\item Measured by meaning, close overlap did not fall (Fact~\ref{fact:wording}). Under cheaper entry it is
nearly pinned, because $\mathcal E$ is small (\cref{prop:volume}).
\end{itemize}
So far AI has acted on science through \emph{entry}. The reading channels remain
predictions. When literature-summarising tools spread, the model predicts
denser topics, fewer new ones, more overshooting by fast-growing topics, and
little or no gain in output per researcher.

\paragraph{Writing AI scales science; reading AI crowds it.} The asymmetry in
\cref{prop:ai} is the paper's sharpest policy implication. Writing speed
multiplies every payoff, including the frontier's. Newcomers spread out, so a
faster pen raises output almost one for one and leaves duplication nearly
unchanged in the long run.
Reading speed improves only existing topics. Newcomers crowd in until those
topics are no better than the frontier, so the improvement is competed away.
This overturns a natural intuition. The worry that writing tools will flood science
with duplicates is misplaced in the long run: free topic choice pins overlap.
The costs of a faster pen are transitional (\cref{app:extensions}) and
concentrated in topics that grow too fast. The hope that
reading tools will raise productivity is also misplaced: free topic choice
dissipates their gains. Tools are worth subsidising only if they make new topics easier
to open, that is, if they raise $F$.
